\documentclass[11pt,a4paper]{article}

\usepackage{lmodern}
\usepackage[T1]{fontenc}
\usepackage[utf8]{inputenc}
\usepackage[margin=2.6cm]{geometry}
\usepackage{amsmath}
\usepackage{amssymb}
\usepackage{booktabs}
\usepackage{microtype}
\usepackage{xcolor}
\usepackage[round]{natbib}
\usepackage[hidelinks]{hyperref}

\newcommand{\code}[1]{\texttt{#1}}
\DeclareMathOperator{\CVaR}{CVaR}
\DeclareMathOperator{\VaR}{VaR}

\title{\code{ifunnel}: A Four-Stage Pipeline for CVaR-Targeted Portfolio
Construction, Measured Stage by Stage}

\author{Thomas Schmelzer\thanks{\url{https://github.com/tschm/funnel}}}

\date{\today}

\begin{document}

\maketitle

\begin{abstract}
\code{ifunnel} builds portfolios by pushing a large fund universe through four
stages: select a diversified subset by minimum spanning tree or by clustering,
estimate moments and generate scenarios, derive a risk target from a benchmark,
and solve a convex programme that maximises expected return subject to a
Conditional Value-at-Risk limit. This paper measures each stage separately on
the $754$-fund Danish universe the repository ships, selecting on $2017$--$2020$
and testing on $2021$--$2024$. The largest effect we can measure anywhere in the
pipeline is not a modelling choice at all: extending the selection window to
overlap the evaluation period lifts the clustering selector's equal-weight
Sharpe ratio from $+0.17$ to $+1.02$, because that selector ranks funds by
realised Sharpe ratio, and nothing in its signature or documentation says the
window must end first. The second largest is the optimisation stage, which on
this window costs between $1.0$ and $1.5$ Sharpe against equal weight on the
same names, in all six model and scenario combinations we ran; two mechanisms
account for it, an objective that extrapolates the training-window mean and a
risk budget taken from the benchmark rather than from the subset --- the
benchmark's $4$-week tail is $1.29$ to $2.00$ times the subset's, and the
optimised portfolios duly come back at two to three times the volatility of
equal weight. Three findings are defects rather than surprises. The Monte Carlo
scenario branch produces essentially Gaussian $4$-week returns (skewness
$+0.10$, excess kurtosis $+0.10$) where the bootstrap branch keeps most of the
historical tail ($-1.05$, $+3.17$), understating the CVaR of the same portfolio
by $33\%$; because the target is \emph{always} bootstrapped, that mismatch
enlarges the risk budget by half rather than cancelling. The Portfolio
Diversification Index is computed from the wrong spectrum and misses its last
term, returning $n-1$ instead of $n$ for $n$ independent assets and NaN for a
perfectly correlated block, and on real data it reverses the advice: the coded
index says diversification improves through five MST passes, the correct one
says it peaks after two. And the CVaR path is unseeded, so an identical call
returns Sharpe ratios spanning up to $0.76$ --- wider than the differences
between the configurations it is used to compare.
\end{abstract}

\section{Introduction}

The portfolio problem that \citet{markowitz1952} posed has a well-known
practical failure mode: the optimiser is a magnifying glass held over the
estimation error in its inputs \citep{michaud1989,best1991}, and the remedy is
usually taken to be either better estimates \citep{ledoit2004,jorion1986} or
fewer decisions \citep{demiguel2009}. Replacing variance with a coherent tail
measure \citep{artzner1999} --- Conditional Value-at-Risk, which
\citet{rockafellar2000} and \citet{krokhmal2002} showed can be optimised as a
linear programme --- changes what is being controlled but not that fact.

\code{ifunnel} is an implementation of a complete pipeline built on that
machinery, aimed at a universe too large to optimise over directly: $754$
Danish mutual funds and ETFs, weekly, from $2013$. Its four stages are

\begin{enumerate}
  \item \textbf{Selection.} Reduce the universe to a few dozen names, either by
        repeatedly taking the leaves of a minimum spanning tree over the
        correlation graph \citep{mantegna1999,onnela2003}, or by hierarchical
        clustering followed by picking the best-performing names in each
        cluster.
  \item \textbf{Scenarios.} Estimate rolling means and Ledoit--Wolf-shrunk
        covariances, and generate $4$-week return scenarios either by Monte
        Carlo from those moments or by bootstrapping historical weeks
        \citep{efron1979}.
  \item \textbf{Targets.} Compute, for each rebalancing period, the CVaR of an
        equal-weight portfolio of the chosen \emph{benchmark}, and use it as
        the risk limit.
  \item \textbf{Optimisation.} Maximise scenario expected return subject to
        that CVaR limit, a budget constraint, linear transaction costs and a
        concentration cap, rebalancing every four weeks.
\end{enumerate}

Each stage is individually defensible and each is a published technique. The
question this paper asks is which of them the result actually depends on. That
question is worth asking of any pipeline, and it is worth asking here in
particular because the four stages differ by two orders of magnitude in code
size, and the two smallest of them turn out to carry the two largest effects we
can measure.

Section~\ref{sec:method} states the pipeline precisely.
Section~\ref{sec:impl} records the implementation choices that determine the
numbers. Section~\ref{sec:results} measures the stages. Where a measurement
turned up something broken rather than merely surprising, we say so and give
the repair.

\section{The pipeline}
\label{sec:method}

Let $R \in \mathbb{R}^{T \times N}$ be weekly returns for $N$ funds.

\subsection{Selection}

\paragraph{Minimum spanning tree.}
Spearman correlations $C$ are mapped to distances
$d_{ij} = \sqrt{2(1 - C_{ij})}$, a complete graph is built on those distances,
and its minimum spanning tree \citep{kruskal1956} is computed. The surviving
subset is the set of \emph{leaves} --- vertices of degree one:
\begin{equation}
  S = \{\, i : \deg_{\mathrm{MST}}(i) = 1 \,\}.
  \label{eq:mst}
\end{equation}
The rationale is that a leaf is a fund the tree could not place at the centre
of a cluster, so leaves are the periphery of the correlation structure. The
operation is applied repeatedly --- $n$ passes, each on the survivors of the
last --- with $n$ a user parameter.

Two diagnostics are returned with the subset: the average correlation among the
survivors, and the Portfolio Diversification Index of \citet{rudin2006},
\begin{equation}
  \mathrm{PDI} = 2 \sum_{k=1}^{n} k \, w_k - 1,
  \qquad
  w_k = \frac{\lambda_k}{\sum_j \lambda_j},
  \label{eq:pdi}
\end{equation}
where $\lambda_1 \geq \cdots \geq \lambda_n$ are the eigenvalues of the
subset's correlation matrix. PDI equals $1$ for a perfectly correlated block
and $n$ for $n$ independent assets, so it reads as an effective number of bets.

\paragraph{Clustering.}
The alternative maps the same correlations to distances, runs complete-linkage
hierarchical clustering into $k$ groups, and keeps the $m$ funds with the
highest Sharpe ratio in each group, giving at most $km$ names. Unlike
\eqref{eq:mst} this stage conditions on realised performance, which is what
Section~\ref{sec:lookahead} is about.

\subsection{Moments and scenarios}

For rebalancing period $p$ the moments are estimated on a rolling window that
ends where period $p$ begins: with $T_{\mathrm{train}} = T - T_{\mathrm{test}}$,
the window is rows $[4p,\; T_{\mathrm{train}} + 4p)$. The covariance is the
biased sample covariance shrunk towards a scaled identity with the
\citet{ledoit2004} intensity, and the mean is the plain window mean.

Scenarios are $4$-week compounded returns, $2^{\text{nd}}$-stage inputs to the
optimiser, generated one of two ways:
\begin{align}
  \text{Monte Carlo:} \quad
    & s^{(p)}_{j} = \prod_{u=1}^{4}\left(1 + z^{(p)}_{j,u}\right) - 1,
      \qquad z^{(p)}_{j,u} \sim \mathcal{N}(\mu_p, \Sigma_p),
      \label{eq:mc}\\
  \text{Bootstrap:} \quad
    & s^{(p)}_{j} = \prod_{u=1}^{4}\left(1 + R_{\tau_{j,u}, \cdot}\right) - 1,
      \qquad \tau_{j,u} \sim \mathrm{Unif}\{4p, \ldots, T_{\mathrm{train}} + 4p\}.
      \label{eq:boot}
\end{align}
The bootstrap draws whole \emph{rows}, so the cross-sectional dependence of a
given week is preserved exactly; the four weeks within a scenario are drawn
independently, so serial dependence is not \citep{politis1994}.

\subsection{Target and optimisation}

The risk limit for period $p$ is the $\CVaR_{\alpha}$, at $\alpha = 5\%$, of an
equal-weight portfolio of the benchmark's assets under \emph{bootstrap}
scenarios of the benchmark --- irrespective of which generator supplies the
optimiser's own scenarios.

Given scenarios $s^{(p)} \in \mathbb{R}^{J \times n}$ with equal probabilities,
previous holdings $x_{\mathrm{old}}$, cash $c$, and the period's target
$\theta_p$, the programme solved is
\begin{equation}
\begin{aligned}
  \max_{x \geq 0} \quad & \bar\mu_p^{\top} x \\
  \text{s.t.} \quad
    & \zeta + \frac{1}{J\alpha} \sum_j \left[ -s^{(p)}_j{}^{\top} x - \zeta \right]_{+}
      \leq \theta_p \cdot V_p, \\
    & \mathbf{1}^{\top} x + \kappa \lVert x - x_{\mathrm{old}} \rVert_1
      = \mathbf{1}^{\top} x_{\mathrm{old}} + c, \\
    & x \leq w_{\max} \, \mathbf{1}^{\top} x,
\end{aligned}
\label{eq:cvar}
\end{equation}
with $\zeta$ the free Value-at-Risk variable, $\kappa$ the linear transaction
cost rate and $V_p$ the portfolio value entering the period. This is the
\citet{rockafellar2000} linearisation, expressed in CVXPY \citep{diamond2016}
with the hinge realised as an auxiliary non-negative vector. The Markowitz
alternative replaces the CVaR constraint by
$x^{\top} \Sigma_p x \leq \theta_p$ with $\theta_p$ from a variance target.

Note the objective. $\bar\mu_p$ is the scenario mean, which is the rolling
window mean --- so the programme maximises extrapolated historical return
subject to a tail constraint. It is the configuration \citet{best1991} and
\citet{schmelzer2013} identify as the fragile one, and it is worth keeping in
mind when reading Section~\ref{sec:backtest}.

\section{Implementation}
\label{sec:impl}

\paragraph{Four weeks is hard-wired.}
\code{n\_iter = 4} appears independently in the moment generator, both scenario
generators and the optimiser loop. The rebalancing frequency is not a
parameter.

\paragraph{Two of the four stages are deterministic; two are not.}
Both selectors are deterministic: the MST is unique here, and the clustering is
complete-linkage on a fixed distance matrix with a Sharpe-ranked pick.
Everything downstream draws from an \emph{unseeded} generator:
\code{backtest} constructs \code{ScenarioGenerator(np.random.default\_rng())}
with no seed. The lifecycle entry point does accept an \code{rng\_seed}, but
treats the value $0$ --- its own default, and the first value most users would
pass --- as ``no seed''.

\paragraph{Costs and caps are fixed inside the call.}
\code{backtest} passes $\kappa = 0.001$, $w_{\max} = 1$, $\alpha = 0.05$ and a
budget of $100$ as literals. They are parameters of the models it calls but not
of the pipeline the user drives.

\paragraph{A solver failure is not reported as one.}
When \eqref{eq:cvar} does not solve to optimality, the rebalancing routine
pickles its inputs to \code{rebalance\_inputs.pkl} in the working directory,
logs, and returns \code{None} --- which the caller immediately unpacks into
four names. The user therefore sees a \code{TypeError} about unpacking
\code{None} rather than the diagnostic that was written for them. In our runs
the CLARABEL path never triggered this, so the observation is from reading, not
from measurement.

\paragraph{The reported Sharpe ratio is a ratio of rounded numbers.}
\code{final\_stats} rounds the annualised return and volatility to two decimals
\emph{before} dividing. For an annual return of $8.4\%$ and volatility of
$12.6\%$ the exact ratio is $0.667$ and the reported one is $0.615$. All Sharpe
ratios in this paper are computed from the portfolio value path instead, and we
note the discrepancy because it is of the same order as some of the differences
the pipeline is used to compare.

\section{Results}
\label{sec:results}

The universe is $754$ Danish funds and ETFs with weekly returns. All selection
uses $2017$-$01$-$04$ to $2020$-$12$-$30$ ($209$ weeks); all evaluation uses
$2021$-$01$-$06$ to $2024$-$07$-$24$ ($186$ weeks, $47$ rebalancing periods).
The benchmark is \code{iShares MSCI World ETF USD Dist}, which returned
$13.26\%$ a year at $14.32\%$ volatility over the test window. Where a number
depends on the unseeded generator we report the mean and range over eight
identical runs.

\subsection{Selection: what the MST passes do}

\begin{table}[t]
\centering
\caption{Repeated MST leaf selection on the selection window. \emph{As
returned} is the average correlation the function reports, which includes the
unit diagonal; \emph{off-diagonal} is the average over distinct pairs. The two
PDI columns are \eqref{eq:pdi} as implemented and as defined.}
\label{tab:mst}
\begin{tabular}{rrrrrr}
\toprule
Pass & Kept & $\bar\rho$ as returned & $\bar\rho$ off-diagonal & PDI as coded & PDI from \eqref{eq:pdi} \\
\midrule
$0$ & $754$ & ---      & $0.4375$ & ---     & ---      \\
$1$ & $313$ & $0.4067$ & $0.4048$ & $2.24$  & $20.43$  \\
$2$ & $147$ & $0.3752$ & $0.3709$ & $2.66$  & $20.91$  \\
$3$ & $77$  & $0.3255$ & $0.3166$ & $3.65$  & $20.05$  \\
$4$ & $41$  & $0.2794$ & $0.2614$ & $5.53$  & $16.75$  \\
$5$ & $19$  & $0.2658$ & $0.2250$ & $7.92$  & $11.23$  \\
$6$ & $11$  & $0.2566$ & $0.1823$ & $6.53$  & $7.87$   \\
\bottomrule
\end{tabular}
\end{table}

Table~\ref{tab:mst} shows the selector doing what it claims. Each pass keeps
roughly half of its input --- $754 \to 313 \to 147 \to 77 \to 41 \to 19 \to 11$
--- and the average pairwise correlation falls monotonically from $0.44$ to
$0.18$. As a mechanism for finding the periphery of a correlation graph it
works.

The diagnostics reported alongside it are another matter, and they disagree
about when to stop.

\paragraph{The average correlation is inflated, increasingly so.}
The value returned is the mean of the full correlation matrix, diagonal
included, which overstates the average pairwise correlation by
$(1 - \bar\rho)/n$. At pass $1$ that is $0.002$ and irrelevant; at pass $6$,
with $n = 11$, it is $0.074$ --- the function reports $0.2566$ where the
average pairwise correlation is $0.1823$. The bias grows exactly as the subset
shrinks, which is to say exactly where the number is being used to argue that
shrinking helped.

\paragraph{The PDI is computed from the wrong spectrum.}
Equation~\eqref{eq:pdi} needs the eigenvalues of the correlation matrix. The
implementation instead fits a PCA \citep{pedregosa2011} to the correlation
matrix \emph{as if it were a data matrix} --- rows treated as observations ---
so it eigendecomposes the covariance of $C$'s rows rather than $C$. The
centring that PCA performs costs one degree of freedom, and the sum in the code
runs to $n-1$ rather than $n$, dropping the last term. Both symptoms are
visible on matrices whose answer is known:

\begin{table}[t]
\centering
\caption{PDI on reference matrices. For $n$ independent assets the index should
be $n$; for a perfectly correlated block, $1$.}
\label{tab:pdi}
\begin{tabular}{rrrrr}
\toprule
& \multicolumn{2}{c}{$C = I_n$ (answer $n$)} & \multicolumn{2}{c}{$C = \mathbf{1}\mathbf{1}^{\top}$ (answer $1$)} \\
\cmidrule(lr){2-3}\cmidrule(lr){4-5}
$n$ & As coded & From \eqref{eq:pdi} & As coded & From \eqref{eq:pdi} \\
\midrule
$5$  & $4.000$  & $5.000$  & NaN & $1.000$ \\
$20$ & $19.000$ & $20.000$ & NaN & $1.000$ \\
$50$ & $49.000$ & $50.000$ & NaN & $1.000$ \\
\bottomrule
\end{tabular}
\end{table}

Table~\ref{tab:pdi} isolates the two faults. The identity case returns exactly
$n-1$, the missing term. The rank-one case returns NaN, because a matrix of
identical rows has zero total variance after centring and the explained-variance
ratio divides by it --- accompanied by a runtime warning, and propagating into
whatever the caller does next.

On real data the consequence is not a small bias but reversed advice. In
Table~\ref{tab:mst} the coded PDI rises through pass $5$ and would tell a user
to keep iterating; the PDI of \eqref{eq:pdi} peaks at pass $2$ ($20.91$) and
falls thereafter, saying the diversification of the subset is being destroyed
from pass $3$ onwards even as the average correlation keeps improving. The
second reading is the coherent one: PDI is bounded above by the number of
assets, and a subset of $11$ funds cannot carry $20$ independent bets however
uncorrelated it is. Repairing \eqref{eq:pdi} is four lines, and it changes the
recommended number of MST passes from ``as many as you like'' to two.

\subsection{Selection: the window is the whole game}
\label{sec:lookahead}

\begin{table}[t]
\centering
\caption{Equal-weight performance over the fixed test window
($2021$-$01$-$06$ to $2024$-$07$-$24$) of subsets chosen by two selection
windows. The second block lets the selection window overlap the evaluation
period, as it does whenever a caller passes the full history.}
\label{tab:lookahead}
\begin{tabular}{llrrrr}
\toprule
Selection window & Selector & $n$ & Return & Vol.\ & Sharpe \\
\midrule
ends $2020$-$12$-$30$ & MST, $4$ passes      & $41$ & $+4.67\%$ & $9.22\%$  & $+0.51$ \\
ends $2020$-$12$-$30$ & clustering $5\times5$ & $25$ & $+1.39\%$ & $8.23\%$  & $+0.17$ \\
\midrule
ends $2024$-$07$-$24$ & MST, $4$ passes      & $32$ & $+3.52\%$ & $10.12\%$ & $+0.35$ \\
ends $2024$-$07$-$24$ & clustering $5\times5$ & $25$ & $+8.12\%$ & $7.97\%$  & $+1.02$ \\
\bottomrule
\end{tabular}
\end{table}

Table~\ref{tab:lookahead} is the most important table in this paper and it is
about an API, not an algorithm. The two selectors are called the same way ---
\code{mst(start\_date, end\_date, n\_mst\_runs)} and
\code{clustering(start\_date, end\_date, n\_clusters, n\_assets)} --- and
neither signature nor docstring says that the window must end before the period
you intend to evaluate.

For the MST selector it does not much matter: it uses only the correlation
matrix, and extending its window into the test period moves the Sharpe ratio
from $+0.51$ to $+0.35$, in the wrong direction. For the clustering selector it
matters enormously: extending the window lifts the equal-weight Sharpe ratio
from $+0.17$ to $+1.02$. The reason is \code{pick\_cluster}, which keeps the
highest-Sharpe funds in each cluster --- so any overlap between the selection
window and the evaluation window is a direct look-ahead on the quantity being
evaluated.

The gap is $0.85$ Sharpe, and it is larger than every difference between models
and scenario generators reported below. A user comparing MST against clustering
with the natural call --- pass the whole history to both, then backtest --- would
conclude that clustering wins by a wide margin. On this data, with the windows
kept honest, it loses by a wide margin under equal weight.

\subsection{Scenarios: which tail}

\begin{table}[t]
\centering
\caption{Equal-weight $4$-week return of the MST subset ($41$ funds), pooled
over all $47$ periods and $2000$ scenarios. \emph{Historical} is
non-overlapping $4$-week compounding of the realised weekly returns. The
Gaussian reference for $\CVaR_{5\%}/\sigma$ is $2.063$.}
\label{tab:scen}
\begin{tabular}{lrrrrr}
\toprule
Source & Mean & SD & Skew & Excess kurt.\ & $\CVaR_{5\%}$ \\
\midrule
Monte Carlo \eqref{eq:mc} & $+0.00499$ & $0.03131$ & $+0.098$ & $+0.097$  & $0.05831$ \\
Bootstrap \eqref{eq:boot} & $+0.00489$ & $0.03290$ & $-1.051$ & $+3.172$  & $0.08664$ \\
Historical                & $+0.00493$ & $0.03702$ & $-2.839$ & $+19.087$ & $0.09137$ \\
\bottomrule
\end{tabular}
\end{table}

Table~\ref{tab:scen} answers a question the \code{scenarios\_type} flag makes
look cosmetic. The Monte Carlo branch samples a Gaussian and compounds four
draws; four is not enough compounding to generate a tail, and the result is a
distribution with skewness $+0.10$ and excess kurtosis $+0.10$. A
$\CVaR$ constraint imposed on it is a variance constraint wearing different
notation. The bootstrap branch, drawing historical weeks, recovers skewness
$-1.05$ and excess kurtosis $+3.17$ --- about $86\%$ of the historical excess
tail, measured as
$(\CVaR_{\mathrm{boot}} - \CVaR_{\mathrm{MC}})/(\CVaR_{\mathrm{hist}} -
\CVaR_{\mathrm{MC}})$. It falls short of history because it compounds four
independently drawn weeks and so destroys the serial dependence that produces
the $-2.84$ skewness of realised $4$-week returns; a block bootstrap
\citep{politis1994} would close part of that gap.

For the same portfolio, then, the Monte Carlo branch reports a $5\%$ CVaR of
$0.0583$ where the bootstrap reports $0.0866$: it understates the tail by
$33\%$.

\paragraph{The mismatch does not cancel.}
It would be a smaller problem if the target and the constraint were computed
the same way, because a uniformly understated risk measure on both sides would
partly cancel. They are not: the target generator always bootstraps, whatever
\code{scenarios\_type} says. So selecting Monte Carlo scenarios leaves a
fat-tailed target governing a thin-tailed constraint, and the effective risk
budget is inflated by the ratio of the two --- $1/0.667 = 1.50$ for this subset.
The configuration that a user would pick to get a faster, smoother scenario
set is the one that quietly loosens the risk limit by half again.

\subsection{The target: the benchmark is the risk budget}

\begin{table}[t]
\centering
\caption{Mean $\CVaR_{5\%}$ of the equal-weight portfolio of each asset set,
averaged over the $47$ rebalancing periods. The target handed to the optimiser
is the benchmark's row.}
\label{tab:target}
\begin{tabular}{lrrr}
\toprule
Asset set & Bootstrap & Monte Carlo & MC / bootstrap \\
\midrule
Benchmark: MSCI World ETF ($1$) & $0.1103$ & $0.0876$ & $0.794$ \\
MST subset ($41$)               & $0.0856$ & $0.0571$ & $0.667$ \\
Clustering subset ($25$)        & $0.0553$ & $0.0450$ & $0.813$ \\
\bottomrule
\end{tabular}
\end{table}

Table~\ref{tab:target} explains a result that is otherwise puzzling: why the
optimised portfolios of Section~\ref{sec:backtest} run at two to three times the
volatility of equal weight on the same names. The risk limit is not a property
of the selected subset. It is the tail of an equal-weight portfolio of the
\emph{benchmark}, and a single world-equity ETF has a much fatter $4$-week tail
($0.1103$) than a basket of $41$ diversified funds ($0.0856$) or $25$
($0.0553$). Choosing that benchmark therefore hands the optimiser a budget
$1.29$ times what the MST subset already runs, and $2.00$ times for the
clustering subset --- and the optimiser, whose objective is expected return,
spends all of it.

This is design rather than defect: the model is a benchmark-relative one, and
targeting the benchmark's CVaR is a coherent thing to want. It is worth stating
plainly because the \code{benchmarks} argument is documented as a list of names
to compare against, its only in-code comment is ``Find Benchmarks' ISIN
codes'', and nothing signals that it is the single most consequential number
in the call.

\subsection{The backtest}
\label{sec:backtest}

\begin{table}[t]
\centering
\caption{Optimised portfolios over the test window, each run eight times with
identical arguments. Ranges are over those eight runs; volatility and drawdown
are means. Rows marked \emph{exact} returned the same value all eight times ---
the Markowitz path never consumes the scenarios it generates. The passive rows
are deterministic.}
\label{tab:backtest}
\small
\begin{tabular}{llrrrr}
\toprule
Subset & Model / scenarios & Return (range) & Vol.\ & Sharpe (range) & Max DD \\
\midrule
\multicolumn{6}{l}{\emph{Passive references}} \\
---     & MSCI World ETF          & $+13.26\%$ & $14.32\%$ & $+0.93$ & $-17.59\%$ \\
all $754$ & equal weight          & $+4.99\%$  & $8.67\%$  & $+0.58$ & $-17.00\%$ \\
MST $41$  & equal weight          & $+4.67\%$  & $9.22\%$  & $+0.51$ & $-16.77\%$ \\
Clust.\ $25$ & equal weight       & $+1.39\%$  & $8.23\%$  & $+0.17$ & $-20.03\%$ \\
\midrule
\multicolumn{6}{l}{\emph{Optimised}} \\
MST $41$ & CVaR / bootstrap & $-10.79\%$ $[-20.7, -5.0]$ & $18.23\%$ & $-0.58$ $[-0.88, -0.30]$ & $-40.69\%$ \\
MST $41$ & CVaR / MC        & $-13.23\%$ $[-17.5, -11.4]$ & $21.37\%$ & $-0.62$ $[-0.84, -0.52]$ & $-47.58\%$ \\
MST $41$ & Markowitz / MC   & $-11.42\%$ (exact) & $20.69\%$ & $-0.55$ (exact) & $-39.56\%$ \\
Clust.\ $25$ & CVaR / bootstrap & $-10.17\%$ $[-18.8, +2.2]$ & $25.63\%$ & $-0.36$ $[-0.64, +0.12]$ & $-57.14\%$ \\
Clust.\ $25$ & CVaR / MC        & $-2.77\%$ $[-8.7, +1.6]$ & $21.80\%$ & $-0.13$ $[-0.38, +0.07]$ & $-41.22\%$ \\
Clust.\ $25$ & Markowitz / MC   & $-0.28\%$ (exact) & $17.95\%$ & $-0.02$ (exact) & $-35.10\%$ \\
\bottomrule
\end{tabular}
\end{table}

Table~\ref{tab:backtest} is the pipeline end to end, and its headline is
uniform: every one of the six optimised configurations loses money over the test
window, and every one trails equal weight on its own subset by a wide margin.
Four things in it are worth separating from that.

\paragraph{The optimisation stage is where the value goes, and the mechanism is
identified.}
Equal weight on the MST subset earns $+4.67\%$; optimising over the same $41$
names earns $-10.8\%$ to $-13.2\%$. The gap is $1.0$ to $1.5$ Sharpe, which is
larger than the difference between the two selectors, larger than the difference
between CVaR and Markowitz, and larger than the difference between the two
scenario generators. We would not read this as evidence about the models: the
objective in \eqref{eq:cvar} maximises the training-window mean, extrapolating
$2017$--$2020$ into a period in which the leadership of this universe changed
completely, and one window cannot separate a bad estimator from a bad four
years. What it does establish is that the stage carrying almost all of the code
is also the stage carrying almost all of the risk of being wrong --- which is
the point \citet{demiguel2009} make about optimisation in general and
\citet{best1991} make about this objective in particular.

\paragraph{The volatility is two to three times equal weight, as
Table~\ref{tab:target} predicts.}
Equal weight on the MST subset runs at $9.22\%$ and on the clustering subset at
$8.23\%$; the optimised versions run at $18\%$ to $26\%$. The CVaR constraint is
not being violated --- it is being \emph{met}, against a target set by a
benchmark whose own tail is $1.29$ to $2.00$ times the subset's. Nothing about
the calls in Table~\ref{tab:backtest} announces that as a risk decision, and it
is the largest one in them.

\paragraph{The selector ranking does not survive the optimiser, and cannot be
measured anyway.}
The clustering subset is the worse of the two under equal weight
($+0.17$ against $+0.51$) and the better of the two after optimisation
($-0.36$, $-0.13$, $-0.02$ against $-0.58$, $-0.62$, $-0.55$). We report the
reversal and decline to interpret it, because the run-to-run spread makes it
unmeasurable: the clustering CVaR/bootstrap configuration alone spans
$-0.64$ to $+0.12$ across eight identical calls.

\paragraph{An identical call is not repeatable --- on the CVaR path.}
Because the scenario generator is constructed unseeded, the four CVaR rows of
Table~\ref{tab:backtest} are eight different portfolios each, with Sharpe
spreads of $0.58$, $0.32$, $0.76$ and $0.45$. Those are wider than every
difference between the six optimised configurations, which span $-0.62$ to
$-0.02$. Reporting a single run --- which is
what the API invites, since \code{backtest} returns one statistics frame --- is
reporting a draw as if it were a mean. Threading a \code{seed} through to
\code{ScenarioGenerator} is a two-line change and would make every number in
this section checkable.

The Markowitz rows, by contrast, reproduce exactly, and for an instructive
reason: \code{backtest} generates the scenario tensor before branching on the
model, and the Markowitz branch passes only the deterministic rolling moments to
its solver. The scenarios are built and discarded. That is why those rows have no
range, and why they run in a fifth of the time.

\section{Limitations}

One universe, one window, one benchmark. The test period
($2021$--$2024$) contains a bond selloff, an equity drawdown and a concentrated
recovery, and any conclusion about which stage helps is conditional on that
sequence. The structural statements --- the clustering selector leaks through
its window, the benchmark sets the budget, Monte Carlo has no tail --- would
survive a different window; the levels in Table~\ref{tab:backtest}, and the sign
of the optimisation's contribution, would not.

The optimised rows use $1000$ scenarios and eight repetitions. Eight is enough
to bound the noise, not to estimate its distribution, and $1000$ scenarios is
at the low end for a $5\%$ tail measure --- the CVaR is then an average over
$50$ scenarios, and its own sampling error contributes to the spread we
attribute to the unseeded generator.

We did not examine the lifecycle module, the glide-path generator or the data
acquisition layer, which together are about a third of the package. We also did
not vary the parameters that \code{backtest} fixes internally: the transaction
cost, the concentration cap and the $5\%$ CVaR level are all left at their
hard-coded values, and the transaction cost in particular ($\kappa = 0.001$ on
a four-weekly rebalance) is a modelling choice worth its own study given the
turnover \eqref{eq:cvar} generates.

Finally, the comparison in Section~\ref{sec:lookahead} attributes the
clustering selector's advantage under a contaminated window to look-ahead. That
is the mechanism, but the size of the effect is specific to this window: a
period in which past winners kept winning would show a smaller gap, and one in
which they reversed could show a negative one.

\section{Conclusion}

Measured stage by stage on this universe and window, the two largest effects in
\code{ifunnel} are not choices between models.

The largest is the selection \emph{window}. Letting it overlap the evaluation
period is worth $+0.85$ Sharpe to the clustering selector, because that selector
ranks by realised Sharpe ratio; the MST selector, which reads only correlations,
shows no such advantage. A user who passes the full history to
\code{clustering(start\_date, end\_date, \ldots)} --- the natural call, and the
one nothing in the signature warns against --- gets a look-ahead larger than any
modelling decision in the package.

The second is the optimisation stage, which on this window costs $1.0$ to $1.5$
Sharpe against equal weight on the same names, in all six configurations. Two
mechanisms are identifiable and both are fixable in principle: the objective
extrapolates the training-window mean, and the risk budget comes from the
benchmark rather than the subset --- which is a coherent design for a
benchmark-relative mandate and an invisible one in a call whose
\code{benchmarks} argument reads as a list of things to plot against.

Three findings are repairs, and we would make them in this order. The CVaR
backtest should take a seed: at present identical calls return Sharpe ratios
spanning up to $0.76$, which makes every comparison the pipeline exists to
support unverifiable. The Portfolio Diversification Index should be computed
from the eigenvalues of the correlation matrix and summed to $n$: as written it
returns $n-1$ for independent assets, NaN for a correlated block, and on this
data recommends five MST passes where the correct index recommends two. And the
CVaR target should be generated the same way as the constraint it governs;
bootstrapping the target while the optimiser sees Gaussian scenarios inflates
the effective risk budget by half.

None of this argues against the design. Reducing $754$ funds to $41$ by graph
structure and then optimising a coherent tail measure over them is a reasonable
thing to build, and the parts implementing published methods --- the CVaR
linearisation, the shrinkage, the MST reduction --- do what they say. The lesson
is about where the leverage sits. Two arguments naming a date range mattered
more than every choice of model and scenario generator in the package.

\appendix

\section{Reproducibility}
\label{app:repro}

All numbers come from the repository's own functions on
\code{all\_etfs\_rets.parquet.gzip} under \code{src/ifunnel/financial\_data/}.
The selection
stage:

\begin{verbatim}
from ifunnel import initialize_bot
from ifunnel.models.mst import minimum_spanning_tree
from ifunnel.models.clustering import cluster, pick_cluster

bot = initialize_bot()
sel = bot.weeklyReturns["2017-01-01":"2020-12-31"]
df = sel.copy()
for _ in range(4):
    mst, df, corr_avg, pdi = minimum_spanning_tree(df)

labels = cluster(sel, 5)
stat = bot.get_stat("2017-01-01", str(sel.index.date[-1]))
clu, _ = pick_cluster(data=sel, stat=stat, ml=labels, n_assets=5)
\end{verbatim}

Table~\ref{tab:pdi} compares the coded index against

\begin{verbatim}
w = np.linalg.eigvalsh(C)[::-1]
w = w / w.sum()
pdi = 2 * sum((i + 1) * w[i] for i in range(len(w))) - 1
\end{verbatim}

Tables~\ref{tab:scen} and~\ref{tab:target} use
\code{MomentGenerator.generate\_sigma\_mu\_for\_test\_periods} followed by
\code{ScenarioGenerator(np.random.default\_rng(0))} and
\code{portfolio\_risk\_target}, with the primal \code{cvar} routine from
\code{cvar\_targets}. Table~\ref{tab:backtest} calls

\begin{verbatim}
bot.backtest(start_train_date="2017-01-01", start_test_date="2020-12-31",
             end_test_date="2024-07-24", subset_of_assets=subset,
             benchmarks=["iShares MSCI World ETF USD Dist"],
             scenarios_type=scen, n_simulations=1000, model=model,
             solver="CLARABEL", lower_bound=0)
\end{verbatim}

eight times per configuration. Because \code{backtest} returns only the rounded
statistics frame, the value path was captured by wrapping
\code{ifunnel.models.main.final\_stats}, and the annualised return, volatility,
Sharpe ratio and drawdown recomputed from it without rounding. That wrapper is
the only modification to the package anywhere in this paper; the seed of the
scenario generator inside \code{backtest} cannot be set from outside, which is
why those rows are ranges.

\bibliographystyle{plainnat}
\bibliography{references}

\end{document}
